\section{Extending a period map from compact types to smooth vectors}

% BEGIN W90G1
\paragraph{W90G1. The source uses the smooth globalization topology.}
Let $G=U(p,q)$ preserve the diagonal signature matrix $J$ of the
given signs, let $K=U(p)\times U(q)$ be its block compact subgroup,
and write $\mathfrak g$ for its real Lie algebra. Here $p+q=n$.
The original question starts with a strongly continuous irreducible
unitary representation $\sigma$ on a Hilbert space $H$ and a linear
map specified on $V=H^{K\text{-fin}}$. These are precisely the
vectors whose $K$-orbit spans a finite-dimensional space.
Harish-Chandra's admissibility theorem and the irreducibility
of the underlying Harish-Chandra module give finite multiplicity of
each compact type and irreducibility of $V$ as a
$(\mathfrak g_{\mathbb C},K)$-module for $H\ne0$.
In particular these vectors are smooth. See
Bernstein--Kr\"otz, \emph{Smooth Fr\'echet globalizations of
Harish-Chandra modules}, \S5.1, Remark 5.1, and Harish-Chandra's
subspace correspondence in \emph{Representations of semisimple
Lie groups I}, Lemmas 33--34 and Theorem 5, pp. 227--228.
Adams, van Leeuwen, Trapa and Vogan state the latter implication for
all real reductive groups in \emph{Unitary representations of real
reductive groups}, Theorem 2.7(1)--(2).
For $U(p,q)$ the scalar centre is compact and acts by one character
in an irreducible unitary representation. The product of this centre
with $SU(p,q)$ maps onto $U(p,q)$ with finite kernel, so the
semisimple result gives the asserted correspondence as well.

The source tests smoothness of $w\in H$ on each product of real
one-parameter subgroups. The products obtained from a real basis of
$\mathfrak g$ give local coordinates at the identity, and translation
gives local coordinates at every group element. Consequently its smooth
subspace is exactly $H^\infty$. For a finite list $S$ of words in
real Lie algebra elements let
\[
 s_S(w)=\sum_{(X_1,\ldots,X_j)\in S}
              \|d\sigma(X_1)\cdots d\sigma(X_j)w\|.
\]
The empty word contributes $\|w\|$. These are the Sobolev seminorms
for the usual smooth-vector topology. All words of bounded length
in one real basis suffice by multilinearity. The smooth vectors of
the given unitary Hilbert representation, with this topology, form
a smooth moderate-growth Fr\'echet globalization of $V$
(Bernstein--Kr\"otz, \S2.4 and Corollary 2.16).

The target $E$ in the source can have arbitrary compact-type
multiplicities. It is complete and Hausdorff, with topology generated
by countably many complex seminorms $p_i$, hence Fr\'echet.
The stated growth estimate and continuity of the orbits give joint
continuity of the action. To see the precise scale in the estimate,
set $s(g)=\max(1,\|g\|_{\rm op},\|g^{-1}\|_{\rm op})$. It is
submultiplicative. As $g^{-1}=Jg^*J$, the source's
$1+\sum_{a,b}|g_{ab}|^2$ is bounded above by $(1+n)s(g)^2$ and
bounds $s(g)$ from below. Thus for each $p_i$ there are a continuous
seminorm $q_i$ and an exponent $N_i$ such that
$p_i(\pi(g)v)\le s(g)^{N_i}q_i(v)$, after absorbing a constant in
$q_i$. This gives local equicontinuity in $g$, as used for joint
continuity. Moreover
\[
 P_i(v)=\sup_{g\in G}s(g)^{-N_i}p_i(\pi(g)v)
\]
lies between $p_i(v)$ and $q_i(v)$, and
$P_i(\pi(h)v)\le s(h)^{N_i}P_i(v)$. The $P_i$ therefore give
the polynomially bounded group-continuous seminorms of
Bernstein--Kr\"otz, Definition 2.6.

Finally every real Lie derivative is a continuous endomorphism of
$E$ by hypothesis. For the orbit map of $v$, its derivative along
$g\exp(tX)$ is $\pi(g)d\pi(X)v$. For a basis of left-invariant
vector fields these derivatives exist everywhere and are continuous.
Iteration gives the same assertion in all orders. In local coordinates
the frame-change coefficients are smooth, so this is smoothness of
the orbit map. The Sobolev seminorms of $E$ obtained by applying
these continuous derivatives induce its given topology, the
zero-order seminorms already generating that topology. Thus $E$ is
a smooth Fr\'echet representation of moderate growth, and in particular
a smooth Schwartz-algebra representation by Proposition 2.20 of the
same paper.
% END W90G1

% BEGIN W90G2
\paragraph{W90G2. Globalization gives the extension and its closed image.}
Let $T_0|_V$ commute with $K$ and the real Lie derivatives.
Complex linearity makes it a $(\mathfrak g_{\mathbb C},K)$-map,
and compact finiteness puts its values in $E^{K\text{-fin}}$.
Bernstein--Kr\"otz, Proposition 11.2, gives the unique continuous
$G$-map $T_\infty:H^\infty\to E$ extending it. This proposition
explicitly allows a smooth Schwartz-algebra target. Continuity
gives, for each target seminorm $p_i$, a bound
$p_i(T_\infty w)\le C_i s_{S_i}(w)$ for a finite list of words.
Indeed the maximum of any finite set of generating seminorms is
bounded by their sum, and their word lists can be concatenated.

Suppose the restriction of $T_0$ to $V$ is nonzero, hence injective
by irreducibility. Let $E_0$ be the closure of $T_\infty(H^\infty)$
in $E$. It is closed and $G$-invariant, so the preceding smoothness
and growth conditions descend to its induced Fr\'echet topology.
The $K$-finite vectors in $H^\infty$ are dense in its smooth
topology. Indeed convolution on $K$ by an approximate identity
converges in every smooth-vector seminorm, and its continuous kernels
can be uniformly approximated by finite sums of matrix coefficients.
It follows that $T_0(V)$ is dense in $E_0$. For an irreducible
compact type $\delta$, its character projection is continuous on
$E_0$. Its image is contained in the closure of $T_0(V[\delta])$,
which is finite-dimensional and closed by W90G1. Conversely every
vector of $T_0(V[\delta])$ is fixed by that projection. Taking finite
sums of types gives $E_0^{K\text{-fin}}=T_0(V)$. This step
handles the possibly infinite multiplicities in the ambient $E$.

Now $E_0$ and $H^\infty$ are smooth moderate-growth Fr\'echet
globalizations of the same Harish-Chandra module.
Bernstein--Kr\"otz, Theorem 10.6 (uniqueness of smooth
globalization), gives a topological $G$-isomorphism between them
extending $T_0|_V$. By density it is $T_\infty$. In particular
its image in $E$ is precisely the closed subspace $E_0$. For every
finite list $S$, continuity of the inverse on $E_0$ gives finitely
many of the ambient seminorms and a constant $C$ with
\[
       s_S(w)\le C\max_{i\in I}p_i(T_\infty w).
\]
The right side is the restriction of one continuous seminorm of $E$,
as the source requires. If $T_0$ vanishes on $V$, the unique smooth
extension is zero and its image is closed. This includes $H=0$.

To obtain the source's linear map on the entire Hilbert space,
choose an algebraic complement to $H^\infty$ in $H$ and extend
by zero on it. Equivariance, uniqueness and both seminorm assertions
have precisely the smooth domain just treated. This proves the full
original $\mathrm{Published.CasselmanWallach}$ for every signature.
For the trivial rank-zero group every vector is smooth and compact
finite, and a nonzero irreducible source is one-dimensional, so the
same assertion follows directly. The result applies to any
subsequent selected real model satisfying its input hypotheses.
% END W90G2
